\chapter*{Conclusion Générale}
\label{chap:conclusion}
\markboth{\MakeUppercase{Conclusion Générale}}{}
\addcontentsline{toc}{chapter}{Conclusion Générale}

Ce projet de fin d’études a été une expérience dense, à la fois sur le plan technique et sur le plan humain. Concevoir Smart Focus, c’était relever un défi assez particulier : construire un système cohérent qui touche à la fois à la vision par ordinateur, au hardware embarqué, au développement backend, à l’application mobile et aux modèles de langage. Autant de domaines qu’on ne maîtrisait pas tous au même niveau au départ, et qui ont demandé des efforts d’apprentissage en parallèle du développement.

Sur le plan des résultats, le système livré est fonctionnel. Le pipeline de vision analyse le comportement de l’utilisateur en temps réel, le score de concentration est calculé de manière continue, les modules pédagogiques — chatbot RAG, quiz, flashcards, planning — sont opérationnels et connectés aux interfaces mobiles. L’intégration IoT avec le Raspberry Pi constitue la couche matérielle qui donne au projet sa dimension concrète et différenciante.

Ce n’est pas un système parfait. La sensibilité aux conditions d’éclairage reste un point à améliorer, et certains modules gagneraient à être testés sur un échantillon plus large d’utilisateurs pour affiner les paramètres. Mais l’architecture est solide, modulaire, et les fondations sont posées pour aller plus loin.

Si on devait continuer ce travail, plusieurs pistes méritent d’être explorées :
\begin{itemize}
    \item améliorer la robustesse du pipeline de vision en conditions d’éclairage difficiles ;
    \item enrichir les fonctionnalités pédagogiques — révision vocale, génération de résumés automatiques ;
    \item affiner les modèles de planification à partir des données réelles collectées sur plusieurs semaines ;
    \item développer une version multi-utilisateurs avec des tableaux de bord comparatifs.
\end{itemize}

Ce projet nous a appris autant par ses succès que par ses blocages. La leçon principale : dans un système aussi hétérogène, la priorité n’est pas la perfection de chaque composant, mais la cohérence de l’ensemble.